\section*{Hoofdstuk \ref{ch:g10:chemical-species} --- Chemische stoffen, natuurlijk en synthetisch}

\begin{solution}{exo:g10:chemical-species:1}
Natuurlijk: vanilline uit een vanillestokje, cafeïne uit koffiebonen.
Synthetisch: aspirine, vanilline uit een fabriek.
\end{solution}

\begin{solution}{exo:g10:chemical-species:2}
Het is dezelfde \omterm{def:g6:pure-substances-and-mixtures:species}{chemische stof}, en een stof heeft dezelfde
eigenschappen, wat haar herkomst ook is.
\end{solution}

\begin{solution}{exo:g10:chemical-species:3}
$R_f = 2.4/6.0 = 0.40$.
\end{solution}

\begin{solution}{exo:g10:chemical-species:4}
Hij koelt de damp en maakt hem weer vloeibaar. Doordat het water aan de
onderkant binnenkomt, vult het de hele mantel (en ontmoet het de heetste
damp als laatste, wat het best koelt).
\end{solution}

\begin{solution}{exo:g10:chemical-species:5}
De stof veel beter \omterm{def:g2:mixing-and-dissolving:dissolve}{oplossen} dan water; niet \omterm{def:g6:solutions-and-solubility:miscible}{mengbaar} zijn met water; zo
weinig mogelijk gevaarlijk zijn en makkelijk te verwijderen.
\end{solution}

\begin{solution}{exo:g10:chemical-species:6}
Ethanol is \omterm{def:g6:solutions-and-solubility:miscible}{mengbaar} met water: er ontstaan geen lagen en er valt niets
te scheiden. Cyclohexaan is niet \omterm{def:g6:solutions-and-solubility:miscible}{mengbaar} met water en lost linalool
goed op.
\end{solution}

\begin{solution}{exo:g10:chemical-species:7}
Water ligt bovenop: \qty{1}{mL} dichloormethaan weegt \qty{1.33}{g},
meer dan \qty{1}{mL} water, dus het zakt naar beneden.
\end{solution}

\begin{solution}{exo:g10:chemical-species:8}
Het smelt te laag (aspirine: \qty{135}{\celsius}) en over een traject:
het monster is onzuiver, of geen aspirine.
\end{solution}

\begin{solution}{exo:g10:chemical-species:9}
Linalool en linalylethanoaat (de twee vlekken liggen op de hoogte van de
twee referenties). Linalylethanoaat is het hoogst gekomen.
\end{solution}

\begin{solution}{exo:g10:chemical-species:10}
\omterm{def:g10:chemical-species:distillation}{Hydrodestillatie}; cyclohexaan aan het destillaat toevoegen en schudden;
de lagen scheiden; het cyclohexaan laten verdampen.
\end{solution}

\begin{solution}{exo:g10:chemical-species:11}
Potlood lost niet op in het \omterm{def:g10:chemical-species:tlc}{eluens}; boven het \omterm{def:g10:chemical-species:tlc}{eluens} worden de stipjes
door het plaatje omhoog meegevoerd in plaats van in de bak op te
lossen.
\end{solution}

\begin{solution}{exo:g10:chemical-species:12}
Cyclohexaan kookt bij \qty{81}{\celsius}, linalool bij \qty{198}{\celsius}:
zacht verwarmen drijft het cyclohexaan uit terwijl het linalool blijft.
\end{solution}

\begin{solution}{exo:g10:chemical-species:13}
Het monster is waarschijnlijk die stof (dezelfde $R_f$ als de natuurlijke
op hetzelfde plaatje) en waarschijnlijk zuiver (één vlek, scherp
smeltpunt).
\end{solution}

\begin{solution}{exo:g10:chemical-species:14}
$0.40 \times 5.5 = \qty{2.2}{cm}$. De andere op $0.75 \times 5.5 =
\qty{4.1}{cm}$ (op de mm nauwkeurig), dus ongeveer \qty{1.9}{cm} boven A.
\end{solution}

\begin{solution}{exo:g10:chemical-species:15}
\qty{200}{mL} kan $1.6 \times 0.200 = \qty{0.32}{g}$ \omterm{def:g2:mixing-and-dissolving:dissolve}{oplossen}: met
\qty{0.20}{g} is het niet verzadigd; er kan nog \qty{0.12}{g} in
\omterm{def:g2:mixing-and-dissolving:dissolve}{oplossen}. De \qty{0.20}{g} die opgelost is, zou met het water verloren
gaan; cyclohexaan, dat linalool veel beter \omterm{def:g2:mixing-and-dissolving:dissolve}{oplost}, haalt er het meeste
van terug.
\end{solution}

\begin{solution}{pb:g10:chemical-species:1}
\textbf{1.} De stoom ervan neemt de geurstoffen uit de bloemen mee naar
de koeler.

\textbf{2.} Hij koelt de damp en maakt hem weer vloeibaar; het koude
water komt aan de onderkant binnen.

\textbf{3.} Niet \omterm{def:g6:solutions-and-solubility:miscible}{mengbaar}; minder dicht dan water (ze drijft).

\textbf{4.} Het laagje is heel dun en een deel van de olie is in het
water opgelost of erin verdeeld: een \omterm{def:g10:chemical-species:extraction}{extractie} haalt het terug.

\textbf{5.} Het lost linalool goed op, is niet \omterm{def:g6:solutions-and-solubility:miscible}{mengbaar} met water, en
kookt bij een lage temperatuur (\qty{81}{\celsius}), dus het is makkelijk
te verwijderen.

\textbf{6.} Nee: ethanol \omterm{def:g2:mixing-and-dissolving:mix}{mengt} met water en er ontstaan geen lagen.

\textbf{7.} Bovenop (\qty{0.78}{g} per mL, minder dan water). Met
dichloormethaan (\qty{1.33}{g} per mL) zou de laag van het \omterm{def:g6:solutions-and-solubility:solute}{oplosmiddel}
onderaan liggen.

\textbf{8.} Geen vlam in de buurt (ontvlambaar); werk in een zuurkast met
handschoenen, en verzamel het afval in plaats van het door de gootsteen
te spoelen (giftig voor in het water levende organismen).

\textbf{9.} Cyclohexaan kookt bij \qty{81}{\celsius}, linalool bij
\qty{198}{\celsius}: zacht verwarmen verwijdert alleen het \omterm{def:g6:solutions-and-solubility:solute}{oplosmiddel}.

\textbf{10.} $R_f(\text{L}) = 2.4/6.0 = 0.40$; $R_f(\text{A}) = 4.5/6.0
= 0.75$.

\textbf{11.} Linalool en linalylethanoaat.

\textbf{12.} Nee: ze bevat minstens twee stoffen; het is een \omterm{def:g6:pure-substances-and-mixtures:mixture}{mengsel}.

\textbf{13.} Zodat alle vlekken onder dezelfde omstandigheden ontwikkeld
worden: alleen dan kun je de hoogtes vergelijken.

\textbf{14.} Haar $R_f$ is gelijk aan die van linalool onder dezelfde
omstandigheden: het is waarschijnlijk linalool, identiek aan de
\omterm{def:g10:chemical-species:natural}{natuurlijke stof}.

\textbf{15.} $1.0 \times 0.90 = \qty{0.90}{g}$.

\textbf{16.} $0.90/100 = \qty{0.90}{\percent}$.

\textbf{17.} $100/0.009 \approx 11\,000$\,g, ongeveer \qty{11}{kg}
bloemen.

\textbf{18.} Ze is uit een plant gewonnen, zonder \omterm{def:g5:irreversible-changes:chemical-change}{chemische verandering}.
Maar haar linalool is dezelfde stof als synthetisch linalool: dezelfde
formule, dezelfde eigenschappen.

\textbf{19.} \qty{0.90}{g} etherische olie uit \qty{100}{g} bloemen.
\end{solution}
